\chapter{Couche 8 — Les données et le réseau}

Deux sujets qu'on range rarement sous « sécurité » et qui en relèvent
pourtant : l'intégrité de l'état, et la confidentialité du transport.

% -----------------------------------------------------------------------------
\mesure{Le cloisonnement entre locataires}{mesure 29 · \fichier{database.py}}

\menace
C'est la faille la plus banale des applications multi-utilisateurs : l'URL
\cmd{/instances/42} affiche la machine numéro 42, et personne ne vérifie
qu'elle appartient bien à celui qui la demande. On change le numéro, on voit
les machines des autres --- et leur mot de passe \cmd{root}.

\parade
Le filtrage par propriétaire est fait \textbf{dans la requête SQL}, pas dans le
gabarit ni dans une condition d'affichage. Une donnée qui n'appartient pas à
l'utilisateur ne quitte jamais la base.

\danslecode
\begin{lstlisting}[language=Python]
requete = requete.where(Rental.user_id == user_id)
...
.where(Rental.user_id == user_id, Instance.status == INSTANCE_RUNNING)
...
.where(Archive.user_id == user_id, Archive.expires_at > utc_now())
\end{lstlisting}

La condition est présente sur les instances, sur les locations \emph{et} sur
les archives --- l'archive du travail d'un autre serait tout aussi
compromettante que sa machine.

\begin{retenir}
Filtrer à l'affichage laisse la donnée arriver jusqu'au gabarit : il suffit
d'une page oubliée, d'une réponse JSON, d'un message d'erreur trop bavard pour
qu'elle sorte. Filtrer dans la requête rend la fuite structurellement
impossible. C'est le test le plus important de la suite, et il porte un nom
explicite : \emph{cloisonnement entre comptes}.
\end{retenir}

\limite
Le cloisonnement repose sur la discipline : chaque nouvelle requête doit porter
la condition. Un modèle de données où la propriété serait imposée par le moteur
(\emph{row-level security} de PostgreSQL) l'appliquerait sans dépendre du
développeur.

% -----------------------------------------------------------------------------
\mesure{Les contraintes d'intégrité}{mesure 30 · \fichier{models.py}}

\menace
Deux machines attribuées au même port sur le même nœud : la seconde ne démarre
pas, ou pire, un utilisateur se connecte à la machine d'un autre. Une
vérification faite en Python ne tient pas sous la concurrence --- deux requêtes
simultanées peuvent lire le même « port libre » avant que l'une n'écrive.

\parade
La règle est tenue par le \textbf{moteur de base de données}, par un
\textbf{index unique partiel} : le couple (nœud, port) est unique, mais
seulement parmi les machines en cours d'exécution.

\danslecode
Le « partiel » est essentiel : sans lui, un port utilisé par une machine
détruite resterait réservé pour toujours, et la plage finirait épuisée. La
contrainte ne s'applique donc qu'aux instances \cmd{running}.

Sept tests vérifient ces règles au niveau du moteur
(\fichier{tests/test\_contraintes.py}) : qu'un port en conflit lève bien une
\cmd{IntegrityError}, que le port est libéré quand l'instance expire, que
l'identifiant et l'adresse de courriel sont uniques, que les clés étrangères
interdisent une location sans utilisateur --- avec
\cmd{PRAGMA foreign\_keys = ON} pour SQLite, qui ne les applique pas par
défaut.

\begin{notion}{Pourquoi le moteur et pas le code}
Une règle tenue par le moteur reste vraie même si une requête concurrente, un
script d'administration ou un futur appel oublient de la vérifier. Une règle
tenue seulement par le code applicatif ne tient que tant que personne ne
contourne ce code. L'énoncé demande explicitement de « tester les contraintes
d'intégrité » : la distinction est exactement celle-là.
\end{notion}

\limite
Les contraintes garantissent la cohérence, pas la justesse métier. Rien
n'empêche d'attribuer un port valide mais déjà pris par un service de l'hôte
--- d'où une vérification complémentaire côté système
(\cmd{port\_is\_free\_on\_host}).

% -----------------------------------------------------------------------------
\mesure{TLS de bout en bout}{mesure 31 · rôle \cmd{reverse\_proxy}}

\menace
Le mot de passe de connexion, le cookie de session et les mots de passe
\cmd{root} circulent entre le navigateur et le serveur. En clair, ils sont
lisibles par quiconque partage le réseau --- un réseau d'école, typiquement.

\parade
nginx termine le TLS devant Gunicorn, avec des protocoles et des suites
choisis plutôt que laissés par défaut.

\danslecode
\begin{lstlisting}
listen 443 ssl http2;
ssl_protocols TLSv1.2 TLSv1.3;
ssl_prefer_server_ciphers on;
ssl_ciphers ECDHE-ECDSA-AES256-GCM-SHA384:ECDHE-ECDSA-CHACHA20-POLY1305:
            ECDHE-ECDSA-AES128-GCM-SHA256;
client_max_body_size 16k;
\end{lstlisting}

Trois décisions :

\begin{itemize}[leftmargin=14pt,itemsep=3pt]
\item \textbf{TLS 1.2 minimum} --- les versions 1.0 et 1.1 sont obsolètes et
      vulnérables ; aucun navigateur en service ne les exige.
\item \textbf{Uniquement des suites ECDHE avec AEAD} : \cmd{ECDHE} assure la
      \emph{confidentialité persistante} (un enregistrement du trafic reste
      indéchiffrable même si la clé du serveur fuite plus tard), et GCM ou
      ChaCha20-Poly1305 chiffrent \emph{et} authentifient.
\item \cmd{client\_max\_body\_size 16k} --- la même borne qu'en applicatif, mais
      appliquée par nginx : une requête démesurée est coupée avant d'atteindre
      Python.
\end{itemize}

\begin{notion}{\cmd{ProxyFix}, ou pourquoi tout le monde avait la même IP}
Derrière un proxy inverse, Flask voit l'adresse du proxy --- \cmd{127.0.0.1}
--- et non celle du visiteur. Le compteur anti-force-brute par adresse
compterait donc toutes les tentatives du monde sur une seule, et bloquerait
tout le monde au premier attaquant. \cmd{ProxyFix} (Werkzeug) lit l'en-tête
\cmd{X-Forwarded-For} posé par nginx et rétablit la vraie adresse. Une mesure
de sécurité mal câblée peut se retourner en déni de service.
\end{notion}

\limite
Le certificat est \textbf{auto-signé} : le navigateur avertit. C'est la seule
alerte que l'analyse dynamique signale encore, et c'est le premier élément des
travaux à court terme (Let's Encrypt). Un certificat auto-signé chiffre
correctement, mais ne prouve pas l'identité du serveur : il ne protège donc pas
d'un attaquant capable de s'interposer.
